% Lösungen Einheit 3
\einheitenkopf*{Einheit 3 von 5 · Exponentialfunktion aufstellen und auswerten}

\erg{28}{a) Start 2022, gesucht 2029 \quad b) Start 2021, gesucht 2026 \quad c) Start 2019, gesucht 2027 \quad d) Start 2023, gesucht 2030 \quad e) $t = 0$ unter 2019, Kreuz bei 2027}
\erg{29}{a) $N_0 = 700$, $q = 1{,}02$ \quad b) $N_0 = 12$, $q = 1{,}15$ \quad c) $N_0 = 48\,000$, $q = 1{,}025$ \quad d) $N_0 = 900$, $q = 1{,}35$ \quad e) $N(t) = 12 \cdot 1{,}15^t$ \quad f) $N(t) = 40 \cdot 0{,}98^t$}
\erg{30}{a) $y = 50\,000 \cdot 0{,}88^x$ \quad b) $2\,500$: Besucher im Eröffnungsjahr (Anfangswert); $1{,}07$: jedes Jahr $7\,\%$ mehr Besucher (Wachstumsfaktor); $t$: Jahre seit der Eröffnung \quad c) $1{,}07 = 1 + 0{,}07$: $100\,\%$ des Vorjahres bleiben, $7\,\%$ kommen dazu. \quad d) $N(t) = 1\,150 \cdot 1{,}028^t$}
\erg{31}{a) $4\,000$ \quad b) $108{,}9$ \quad c) $2\,205$ \quad d) $450$ \quad e) $2\,315{,}25$ \quad f) $t = 8$: $N(8) \approx 3\,984$ Einwohner \quad g) $N(7) \approx 1\,590{,}35\,€$ \quad h) $N(5) = 2\,400 \cdot 0{,}93^5 \approx 1\,669{,}65\,€ < 1\,700\,€$ – die Behauptung stimmt. \quad i) A: $18\,000 \cdot 1{,}02^{12} \approx 22\,828$; B unterschätzt um etwa $508$ Einwohner, weil die $2\,\%$ jedes Jahr vom größeren Wert genommen werden.}
\erg{32}{a) nach $4$ Wochen: $N(4) \approx 2\,491$ Pflanzen \quad b) nach $7$ Tagen: $900 \cdot 0{,}85^7 \approx 288{,}5\,\text{kg}$ \quad c) $N(12) = 3\,500 \cdot 1{,}25^{12} \approx 50\,932 > 50\,000$ – die Behauptung stimmt.}
\erg{33}{a) Fehler: Ole hat mit $8$ malgenommen statt mit $8$ zu potenzieren. $N(8) = 600 \cdot 1{,}05^8 \approx 886{,}47$ \quad b) $750 \cdot 1{,}1^6 \approx 1\,328{,}67$ \quad c) Fehler: Ella hat die Jahreszahl eingesetzt statt der Zahl der Schritte; $t = 2031 - 2025 = 6$: $N(6) = 1\,800 \cdot 1{,}02^6 \approx 2\,027{,}09$ \quad d) $t = 7$: $2\,600 \cdot 1{,}015^7 \approx 2\,885{,}60$}
\erg{34}{a) Jedes Jahr wird noch einmal mit $1{,}1$ malgenommen; $t$-mal malnehmen ist die Potenz $1{,}1^t$. Mit $1{,}1 \cdot t$ würde man nur einmal malnehmen. \quad b) Ja: Die $2\,\%$ werden jedes Jahr vom größeren Wert genommen, also wächst es um mehr als $12 \cdot 2\,\%$ (genau $1{,}02^{12} \approx 1{,}268$, etwa $26{,}8\,\%$).}
\erg{35}{a) $N(t) = 250 \cdot 0{,}82^t$ \quad b) $N(3) \approx 137{,}8\,\text{mg}$ \quad c) nach $6$ Stunden ($N(5) \approx 92{,}7\,\text{mg}$, $N(6) \approx 76{,}0\,\text{mg}$)}
